\section{James Gosling}
Criador do Java: o Computer History Museum o fez Fellow em 2019 \enquote{pela
concepção, projeto e implementação da linguagem de programação Java}
(fonte: \citeobra{bio:gosling-chm}).

\begin{itemize}
\obra{paper:java-language-environment}{O white paper de 1996, com Henry
  McGilton: os objetivos de design (simples, orientada a objetos, robusta,
  portável) e o que foi deixado de fora do C e do C++. A melhor resposta a
  \enquote{por que o Java é assim?}.}
\obra{paper:feel-of-java}{Gosling, em primeira pessoa, sobre as decisões de
  design e por que o Java é uma linguagem \enquote{blue collar}.}
\obra{doc:jls25}{A especificação da linguagem, que ainda leva o nome dele
  entre os autores (com Joy, Steele, Bracha, Buckley, Smith e Bierman).}
\obra{livro:arnold2005}{O livro-texto da linguagem na época do J2SE~5.0
  (genéricos, enums, anotações), com Ken Arnold e David Holmes. Hoje vale
  como documento histórico.}
\obra{artigo:gosling-oral-history}{História oral do Computer History Museum
  (parte~1, transcrição gratuita): formação, CMU, Emacs e o caminho até o
  Green Project.}
\obra{video:gosling-foojay-30}{Podcast dos 30 anos do Java: das origens em
  controle de dispositivos aos servidores, e o que ele mudaria hoje.}
\end{itemize}

\section{Tim Lindholm}
Ex-Distinguished Engineer da Sun, contribuiu com a linguagem e foi o
arquiteto sênior da JVM; antes, trabalhou com VMs de Prolog (fonte:
\citeobra{bio:lindholm-informit}).

\begin{itemize}
\obra{doc:jvms25}{A especificação da JVM, que começou com ele e Frank
  Yellin: formato \texttt{.class}, verificação e cada instrução. É a
  referência dos capítulos de bytecode e de JVM da Parte~\ref{parte:java-jvm}.}
\obra{livro:lindholm2014}{A última edição impressa da JVMS (Java~SE~8).
  O item~1.1, \enquote{A Bit of History}, conta a origem da JVM.}
\obra{livro:lindholm1999}{A segunda edição, da era Sun. Serve para comparar
  com a JVM de hoje.}
\obra*{artigo:lindholm-register}{Entrevista de 2026: a reescrita às pressas
  das threads poucos dias antes da estreia do Java, e detalhes que o
  documentário não mostra.}
\end{itemize}
